\clearpage
\section*{Text Processing}
\phantomsection
\addcontentsline{toc}{section}{Text Processing}

% ============================================================================= 
% NORMALIZATION
% ============================================================================= 
\subsection*{Surface Level Normalization Definition}
\phantomsection
\addcontentsline{toc}{subsection}{Surface Level Normalization Definition}

Before tokenizing and any further text processing, the text undergoes a 
three-stage cleaning process to ensure variations in formatting do not create 
duplicate tokens in our index. 
\begin{enumerate}
    \item   The \texttt{.lower()} function converts every character in 
            the documents to lowercase. 
    \item   We strip out any punctuation, including periods, commas, quotes, 
            brackets, and special symbols using a regular expression. Specifically, 
            \texttt{re.sub(r'[\textasciicircum \textbackslash w\textbackslash s]',
            '')}, which tells Python to find any character that is $not$ a word 
            character (numbers included) or a whitespace character, and replace 
            it with nothing. 
    \item   We again use a regular expression to target any sequence of one 
            or more spaces to collapse them to a single, clean space. 
            Additionally, we apply the \texttt{.strip()} function to this 
            regular expression to trim any lingering spaces at the beginning 
            and end of a document. 
\end{enumerate}


% =============================================================================
% TOKENIZER
% =============================================================================
\subsection*{Tokenizer Definition}
\phantomsection
\addcontentsline{toc}{subsection}{Tokenizer Definition}

Now that the text has been normalized, we call \texttt{.split()} to break 
the strings apart at every single space. This function call results in a 
structured list of tokens. 

% ============================================================================= 
% STOPPING
% ============================================================================= 
\subsection*{Stopping Definition}
\phantomsection
\addcontentsline{toc}{subsection}{Stopping Definition}

At the same time the \texttt{.split()} function is tokenizing the text, we 
immediately check the current token against the \texttt{ENGLISH\_STOP\_WORDS}
list provided by the \texttt{scikit-learn} library. If the token is a highly 
common word, (i.e., "the", "is", "at", etc.), the token is instantly discarded. 
This results in a final tokens list full of meaningful words. 

% ============================================================================= 
% STEMMING/LEMMATIZATION
% ============================================================================= 
\subsection*{Stemming/Lemmatizing Definition}
\phantomsection
\addcontentsline{toc}{subsection}{Stemming/Lemmatizing Definition}

Finally, the text processing pipeline ends by looping through the tokens 
and it applies \texttt{NLTK}'s \texttt{PorterStemmer()}, which removes 
common English suffixes and prefixes. This process shrinks our vocabulary 
drastically. 